\section{Conclusion}

In this paper, we proposed DVCL, a dual-view contrastive learning framework for robust heterogeneous graph neural networks. DVCL constructs a purified topology view from meta-path-induced semantic structures and a feature-induced view from target-node similarity. The topology view preserves heterogeneous semantic information through meta-path construction, edge confidence estimation, semantic attention, and two-stage filtering, while the feature view provides complementary evidence that is less dependent on attacked graph structure. A symmetric cross-view contrastive loss aligns the two views and improves representation consistency under structural perturbations.

The first version of the experimental section defines the evaluation protocol for clean graphs, global structural attacks, and target attacks on heterogeneous graph benchmarks. After filling in the measured results, the final analysis will quantify the robustness gains of DVCL, the contribution of each component, and the sensitivity to key hyperparameters. Future work may extend DVCL to dynamic heterogeneous graphs and investigate more adaptive attacks that jointly manipulate topology and node attributes.

% ========== References ==========

